\subsection{等式的性质}

从现在起，我们只考虑等号理论。设 \(\mathfrak{G}\) 为这样的理论，并令 \(\mathfrak{G}_0\) 为这样的理论：其符号为 \(\mathfrak{G}\) 的符号，且其唯一公理由方案S1至S7提供。理论 \(\mathfrak{G}_0\) 弱于 \(\mathfrak{G}\) （§2，第4节）且没有常量。以下三个定理是 \(\mathfrak{G}_0\) 中的定理 .

定理1. x = x。

设 \(S\) 表示关系 \(x = x\)（在 \(\mathcal{C}_0\) 中）。根据C27（§4，第1号），对于每一个关系 \(R\)（在 \(\mathcal{C}_0\) 中）， \((\forall x)(R \leftrightarrow R)\) 是 \(\mathcal{C}_0\) 中的定理 ，因此，根据S7， \(\tau_x(R) = \tau_x(R)\) ，也就是说 \((\tau_x(R)|x)S\) ，是 \(\mathcal{C}_0\) 中的一个定理。取 \(R\) 为关系“非 \(S\) ''，并考虑C26（§4，第1节），我们看到 \((\forall x)S\) 是 \(\mathcal{C}_0\) 中的定理 。根据C30（§4，第3节）， \(S\) 因而是 \(\mathcal{C}_0\) 中的一个定理 .

关系 \((\forall x)(x = x)\) 也是 \(\mathfrak{C}_0\) 中的一个定理；并且如果 \(T\) 是 \(\mathfrak{C}_0\) 中的一项，那么 \(T = T\) 是 \(\mathfrak{C}_0\) 中的定理 （参见§4，第3号）。可以将后面的定理以同样的方式转化为不含任何字母的定理，或转化为元数学判据。从现在起，我们不再明确地进行这些转化，但我们将经常隐式地加以利用。

定理2。 \((x = y) \Longleftrightarrow (y = x)\) .

假设关系 \(x = y\) 成立。由S6，关系

\[
(x = y) \Longrightarrow ((x | y) (y = x) \Longleftrightarrow (y | y) (y = x)),
\]

即

\[
(x = y) \Longrightarrow ((x = x) \Longleftrightarrow (y = x)),
\]

成立。因此 \((x = x) \Longleftrightarrow (y = x)\) 成立。由定理1可知 y = x 成立，定理得证。

定理3。 \(((x = y)\) 和 \((y = z))\Longrightarrow (x = z)\) .

让我们将假设 \(x = y\) , \(y = z\) 附加到 \(\mathfrak{G}_0\) 的公理。由S6，关系 \((x = y) \Longrightarrow ((x = z) \Longleftrightarrow (y = z))\) 成立。因此，

\[
(x = z) \Longleftrightarrow (y = z),
\]

，从而 \(x = z\) 成立。

C44。设 \(\pmb{x}\) 为一个字母，并设 \(\pmb{T}, \pmb{U}, V\{\pmb{x}\}\) 为 \(\mathfrak{C}_0\) 中的项。则关系 \((\pmb{T} = \pmb{U}) \Longrightarrow (V\{\pmb{T}\} = V\{\pmb{U}\})\) 是 \(\mathfrak{C}_0\) 中的定理 .

设 y 和 z 为两个不同的字母，它们与 x 不同，也与 T、U、V 中出现的字母不同。附加假设 y = z。于是，由 S6，

\[
((y | z) (V \{y \} = V \{z \})) \Longleftrightarrow (V \{y \} = V \{z \}),
\]

也就是说 \((V\{y\} = V\{y\}) \Longleftrightarrow (V\{y\} = V\{z\})\) ，成立。现在， \(V\{y\} = V\{y\}\) 由定理 1 成立；因此 \(V\{y\} = V\{z\}\) 为真。

由此一切可得 \((y = z) \Longrightarrow (V\{y\} = V\{z\})\) 是 \(\mathfrak{G}_0\) 中的定理 ，记为 \(A\) 。但 \((T|y)(U|z)A\) 正是

\[
(T = U) \Longrightarrow (V \{T \} = V \{U \}).
\]

形如 \(T = U\) 的关系，其中 \(T\) 和 \(U\) 是 \(\mathfrak{G}\) 中的项，称为方程；（在 \(\mathfrak{G}\) 中）关系 \(T = U\) ——视为字母 \(x\) 中的方程——的解，因此（§2，第2节）是一项 \(V\)（在 \(\mathfrak{G}\) 中），使得 \(T\{V\} = U\{V\}\) 是 \(\mathfrak{G}\) 中的定理 .

设 \(\pmb{T}\) 和 \(\pmb{U}\) 为 \(\mathfrak{C}\) 中的两个项，并且令 \(x_1, x_2, \ldots, x_n\) 为出现在 \(\pmb{T}\) 中但不在 \(\pmb{U}\) 中的字母。如果关系

\[
(\exists x _ {1}) \dots (\exists x _ {n}) (T = U)
\]

是 \(\mathfrak{C}\) 中的定理 时，我们说 \(U\) 可以写成 \(T\) 的形式（在 \(\mathfrak{C}\) 中）。设 \(R\) 是 \(\mathfrak{C}\) 中的一个关系，并令 \(y\) 为一个字母。设 \(V\) 是（在 \(\mathfrak{C}\) 中） \(R\) ——视为 \(y\) 中的一个关系——的一个解。如果（在 \(\mathfrak{C}\) 中） \(R\) ——视为 \(y\) 中的一个关系——的每个解都可以写成 \(V\) 的形式，那么 \(V\) 被称为 \(R\) 的完整解（或通解）（在 \(\mathfrak{C}\) 中）.

